% Options for packages loaded elsewhere
% Options for packages loaded elsewhere
\PassOptionsToPackage{unicode}{hyperref}
\PassOptionsToPackage{hyphens}{url}
\PassOptionsToPackage{dvipsnames,svgnames,x11names}{xcolor}
%
\documentclass[
]{pnas-new}
\usepackage{xcolor}
\usepackage{amsmath,amssymb}
\setcounter{secnumdepth}{-\maxdimen} % remove section numbering
\usepackage{iftex}
\ifPDFTeX
  \usepackage[T1]{fontenc}
  \usepackage[utf8]{inputenc}
  \usepackage{textcomp} % provide euro and other symbols
\else % if luatex or xetex
  \usepackage{unicode-math} % this also loads fontspec
  \defaultfontfeatures{Scale=MatchLowercase}
  \defaultfontfeatures[\rmfamily]{Ligatures=TeX,Scale=1}
\fi
\usepackage{lmodern}
\ifPDFTeX\else
  % xetex/luatex font selection
\fi
% Use upquote if available, for straight quotes in verbatim environments
\IfFileExists{upquote.sty}{\usepackage{upquote}}{}
\IfFileExists{microtype.sty}{% use microtype if available
  \usepackage[]{microtype}
  \UseMicrotypeSet[protrusion]{basicmath} % disable protrusion for tt fonts
}{}
\makeatletter
\@ifundefined{KOMAClassName}{% if non-KOMA class
  \IfFileExists{parskip.sty}{%
    \usepackage{parskip}
  }{% else
    \setlength{\parindent}{0pt}
    \setlength{\parskip}{6pt plus 2pt minus 1pt}}
}{% if KOMA class
  \KOMAoptions{parskip=half}}
\makeatother
% Make \paragraph and \subparagraph free-standing
\makeatletter
\ifx\paragraph\undefined\else
  \let\oldparagraph\paragraph
  \renewcommand{\paragraph}{
    \@ifstar
      \xxxParagraphStar
      \xxxParagraphNoStar
  }
  \newcommand{\xxxParagraphStar}[1]{\oldparagraph*{#1}\mbox{}}
  \newcommand{\xxxParagraphNoStar}[1]{\oldparagraph{#1}\mbox{}}
\fi
\ifx\subparagraph\undefined\else
  \let\oldsubparagraph\subparagraph
  \renewcommand{\subparagraph}{
    \@ifstar
      \xxxSubParagraphStar
      \xxxSubParagraphNoStar
  }
  \newcommand{\xxxSubParagraphStar}[1]{\oldsubparagraph*{#1}\mbox{}}
  \newcommand{\xxxSubParagraphNoStar}[1]{\oldsubparagraph{#1}\mbox{}}
\fi
\makeatother


\usepackage{longtable,booktabs,array}
\usepackage{calc} % for calculating minipage widths
% Correct order of tables after \paragraph or \subparagraph
\usepackage{etoolbox}
\makeatletter
\patchcmd\longtable{\par}{\if@noskipsec\mbox{}\fi\par}{}{}
\makeatother
% Allow footnotes in longtable head/foot
\IfFileExists{footnotehyper.sty}{\usepackage{footnotehyper}}{\usepackage{footnote}}
\makesavenoteenv{longtable}
\usepackage{graphicx}
\makeatletter
\newsavebox\pandoc@box
\newcommand*\pandocbounded[1]{% scales image to fit in text height/width
  \sbox\pandoc@box{#1}%
  \Gscale@div\@tempa{\textheight}{\dimexpr\ht\pandoc@box+\dp\pandoc@box\relax}%
  \Gscale@div\@tempb{\linewidth}{\wd\pandoc@box}%
  \ifdim\@tempb\p@<\@tempa\p@\let\@tempa\@tempb\fi% select the smaller of both
  \ifdim\@tempa\p@<\p@\scalebox{\@tempa}{\usebox\pandoc@box}%
  \else\usebox{\pandoc@box}%
  \fi%
}
% Set default figure placement to htbp
\def\fps@figure{htbp}
\makeatother





\setlength{\emergencystretch}{3em} % prevent overfull lines

\providecommand{\tightlist}{%
  \setlength{\itemsep}{0pt}\setlength{\parskip}{0pt}}



 


\makeatletter
\let\oldlt\longtable
\let\endoldlt\endlongtable

\def\longtable{\@ifnextchar[\longtable@i \longtable@ii}
\def\longtable@i[#1]{%
  % Temporarily fool longtable into thinking we're in onecolumn mode
  \@twocolumnfalse
  \oldlt[#1]%
}
\def\longtable@ii{%
  % Temporarily fool longtable into thinking we're in onecolumn mode  
  \@twocolumnfalse
  \oldlt
}
\def\endlongtable{%
  \endoldlt
  % Restore twocolumn flag
  \@twocolumntrue
}
\makeatother
\makeatletter
\@ifpackageloaded{caption}{}{\usepackage{caption}}
\AtBeginDocument{%
\ifdefined\contentsname
  \renewcommand*\contentsname{Table of contents}
\else
  \newcommand\contentsname{Table of contents}
\fi
\ifdefined\listfigurename
  \renewcommand*\listfigurename{List of Figures}
\else
  \newcommand\listfigurename{List of Figures}
\fi
\ifdefined\listtablename
  \renewcommand*\listtablename{List of Tables}
\else
  \newcommand\listtablename{List of Tables}
\fi
\ifdefined\figurename
  \renewcommand*\figurename{Figure}
\else
  \newcommand\figurename{Figure}
\fi
\ifdefined\tablename
  \renewcommand*\tablename{Table}
\else
  \newcommand\tablename{Table}
\fi
}
\@ifpackageloaded{float}{}{\usepackage{float}}
\floatstyle{ruled}
\@ifundefined{c@chapter}{\newfloat{codelisting}{h}{lop}}{\newfloat{codelisting}{h}{lop}[chapter]}
\floatname{codelisting}{Listing}
\newcommand*\listoflistings{\listof{codelisting}{List of Listings}}
\makeatother
\makeatletter
\makeatother
\makeatletter
\@ifpackageloaded{caption}{}{\usepackage{caption}}
\@ifpackageloaded{subcaption}{}{\usepackage{subcaption}}
\makeatother
\usepackage{bookmark}
\IfFileExists{xurl.sty}{\usepackage{xurl}}{} % add URL line breaks if available
\urlstyle{same}
\hypersetup{
  pdftitle={Variability and acceleration in the growth of children's early vocabulary},
  pdfauthor={Michael C. Frank},
  pdfkeywords={Early word learning; Bayesian data analysis; Large
language models},
  colorlinks=true,
  linkcolor={blue},
  filecolor={Maroon},
  citecolor={Blue},
  urlcolor={Blue},
  pdfcreator={LaTeX via pandoc}}


\templatetype{pnasresearcharticle}

\title{Variability and acceleration in the growth of children's early
vocabulary}

\newcommand{\equalcont}{1}

\newcommand{\correspond}{1}

\author[aff-1%
%
%
]{Michael C. Frank}

\affil[aff-1]{Department of Psychology, Stanford University}







\keywords{Early word learning; Bayesian data analysis; Large language
models}

\begin{document}
\maketitle

\begin{abstract}
Children's early word learning is alternatively portrayed as either a
process of gradual accumulation or as an inferential process that
results in rapid mapping. How can these two views be rectified? Here we
present a statistical model that describes both the acceleration of
children's vocabulary growth and its remarkable variability across
children. This model shows how changes in children's learning efficiency
mean that younger children accumulate knowledge slowly while older
children make rapid inferences. Further, the model can be used to
estimate how much variation between children is due to differences in
home language environment, a key policy-relevant question. Estimates
from four longitudinal datasets converge with a recent meta-analysis in
suggesting XYZ\% of total variance. Finally, we show that, in contrast
to children, large language models are not well-described by this model:
they show standard accumulator dynamics and very limited variability
across a number of dimensions. These results suggest \ldots{}
\end{abstract}

\dates{This manuscript was compiled on \today}
\doi{\url{www.pnas.org/cgi/doi/10.1073/pnas.XXXXXXXXXX}}

\thispagestyle{firststyle}
\ifthenelse{\boolean{shortarticle}}{\ifthenelse{\boolean{singlecolumn}}{\abscontentformatted}{\abscontent}}{}

Children's early language learning is often portrayed as a process of
accumulation and gradual learning, with statistical learning, mutual
exclusivity and input viewpoints highlighting the accumulation nature.
In contrast, studies of processing and word mapping highlight the
increases in efficiency that happen over the same period. How can these
two views be rectified? One clue comes from examining the two key
signatures of early vocabulary growth: variability and acceleration.
What is the shape of this acceleration? And what is the source of
variability, from children's input or otherwise?? We present a
quantitative model that unify these views while explaining both the
variability and acceleration of children's early language growth. The
resulting model provides a good quantitative explanation for longitude
vocabulary growth across large population of children. It also shows
that by their increasing efficiency, and their ability of children,
learning fundamentally different ways than large language, models, which
show classic accumulator dynamics.

Children's early vocabulary grows very rapidly during early childhood,
with large variation between children. The nature of this growth has
theoretical, practical, and methodological importance.

From a theoretical perspective, growth mechanisms shed light on
children's fundamental learning mechanisms and how they change across
childhood. From a practical perspective, variation in growth is an
important precursor --- and predictor --- of academic achievement. From
a methodological perspective, vocabulary also affords a unique
psychometric opportunity to study human variation because, unlike nearly
all important human psychological metrics, there is a true zero point -
a time before which children know no words. Thus, absolute levels of
variability can be calculated.

One dominant viewpoint is that this growth can be characterized as a
process of accumulation (McMurray 2007; Mitchell \& McMurray, 2009;
Hidaka, 2013; Mollica \& Piantadosi, 2017; Kachergis et al., 2021). This
pure accumulator view connects with a viewpoint in which the quantity of
language that children hear has a strong causal role in the speed of
their learning (Fernald papers; Bergeleson 1001). It also connects with
the idea that large language models -- whose learning scales in
predictable, accumulator-like ways (Kaplan, Chinchilla, scaling laws) --
might be good models for children's learning.

But important parts of this proposal have not been fleshed out. For
example, how much variation is due to input quantity vs.~other sources?
Meta-analysis of this relationship suggests it's there but not that big
(Coffey \& Snedeker 2026). Could be quality as well, of course.

Furthermore, how does the accumulation process relate to the rapid
growth of vocabulary? On some models, accumulation by itself can produce
an exponential ``vocabulary explosion'' while the child's own learning
abilities remain constant. On other accounts, however, the child's own
abilities change over development. These changes could be due to the
accumulation of language allowing bootstrapping (e.g., via mutual
excusivity or the shape bias), they could be through faster processing
(Fernald rich get richer), or they could be via general maturation
(e.g., better memory, faster processing). Or all of these.

We address these questions by creating Bayesian data analysis models
that connect proposals about the nature of the accumulation process with
large-scale data about the longitudinal growth of children's vocabulary
over time. Our key innovations are 1) that we model population
variability between children, not just the central tendency and 2) that
we directly connect quantitative estimates of language exposure to the
rate of vocabulary growth.

Critically, we make use of new data resources, first connecting between
vocabulary growth and data-driven estimates of the general range of
language input quantity available to children (study 1), then to
estimates of specific children's language input from BabyView and
SEEDLings datasets (study 2), and finally to estimates of children's
processing speed from Peekbank (study 3).

Overall, our analysis supports a view in which children's vocabulary
growth is not a pure process of accumulation but rather one that
accelerates rapidly over time. Further, although variation in input rate
accounts for significant variation in learning speed, it is limited to
approximately 10\% of total variation across datasets.

Finally, we explore an application of the same framework to large
language models. Making use of GPT-2 models trained on human
developmental data, we show that -- in contrast to children -- these
models do appear to function as pure accumulators. This analysis exposes
a critical disanalogy between LLMs and human learners and suggests a
potential route for understanding the vast ``data gap'' between children
and large language models (Frank, 2023).

Here we ask whether these two signatures of vocabulary growth -
variation and acceleration - are consistent with pure accumulator
models, which are prevalent in the word learning literature and which
characterize the scaling dynamics of large language models.

\section{Model}\label{model}

We start with a simple accumulator model that merges ideas from item
response theory and learning theory (McMurray, Kachergis). We begin with
the Rasch model.

Child \(i\) produces word \(j\) has probability (\(\eta_ij\))

\[P(w_{i,j}=1 \mid \theta_i, \delta_j) =\frac{\exp(\theta_i - \delta_j)}{1 + \exp(\theta_i - \delta_j)}\]
This is an IRT model -- specifically a Rasch model -- where \(\theta_i\)
is child \(i\)'s ability and \(\delta_j\) is word \(j\)'s difficulty.
Following (\textbf{kachergis2021?}), we consider a simple accumulator
model (``pure accumulator''), specifying that \(\theta_i\) is not a
single parameter for a child but instead a function of accumulation over
time. In the simplest form:

\[\theta_i  = log(r_i) + log(t)\]

\noindent where \(r_i\) is the child's rate of language innput (in
tokens per month) and \(t\) is the child's age in months. In the
exponential, this formulation produces a simple rate X time computation.
In other words, ability increases over time. This form instantiates the
basic hypothesis that the only thing that changes over time is more
input, and the only source of variation between children is variation in
rate.

This model does not have any way to capture variation across children.
Thus, we augment this model by adding this model in two ways. First, we
add variation related to accumulation rate:

\[\theta_i  = log(\alpha_i) + log(r_i) + log(t)\]

In this formulation,

Second, we add a temporal acceleration component (an exponent):

\[\theta_i  = log(\alpha_i) + log(r_i) + \kappa_i * log(t)\]

When exponentiated, this turns into
\[\alpha_i * r_i  * t \exp{\kappa_i}\].

This final model instantiates the idea that children vary in their base
level of variation (\(\alpha\)) and their acceleration rate
(\(\kappa\)).

As we will show, such a model fits a range of data vastly better than
models without these terms.

\subsection{Model comparison}\label{model-comparison}

We used the standard reduction of the Rasch model to a mixed effects
regression model (\textbf{jss\_paper?}) to fit the set of models
described above.

We fit these models to data on children's language production. Our
primary data source is the MacArthur-Bates Communicative Development
Inventories {[}MB-CDI;
(\textbf{fenson1994?});(\textbf{marchman2021?}){]}, a popular parent
report instrument in which parents are asked to indicate which words
their child produces and understand. Here we focus exclusively on
production, which is used with older children and which tends to be a
more reliable measure {[}(\textbf{frank2021?});others{]}. We make use of
data from Wordbank (\textbf{frank2017?}), an open repository of data
from the MB-CDI. Since our efficiency variation model is not
identifiable in cross-sectional data, we focus on those datasets that
include substantial amounts of longitudinal data, which come from
American English (N=XYZ), Norwegian (N=XYZ), Quebec French (N=XYZ), and
Japanese (N=XYZ).

The efficiency change was the best fitting model across all four
datasets (as well as two smaller ones; see SI), with AIC differences of
XYZ--XYZ. Figure 1 shows the four primary model variants: pure
accumulator models and models that include population-level efficiency
change predict no variation (A, B). Allowing both efficiency (intercept)
and efficiency change (slope) to vary across individuals results in
substantially better fit.

Power-law growth models (in which growth in ability is a function of the
logarithm of age) showed statistically better fits than exponential
growth models (in which growth in ability is a linear function of age).
Though these models were better from a statistical perspective
(\(\delta\)AIC XYZ -- XYZ), differences were small within the age range
we examined (see Supplemental Table XYZ and Figure XYZ for all model
fits).

\pandocbounded{\includegraphics[keepaspectratio]{../figs/longitudinal/glmer_ladder_main.png}}
The best-fitting modified accumulator model included large developmental
changes in efficiency across all languages (\(\beta = XYZ - XYZ\)). This
finding suggests that we can robustly reject pure accumulation models of
early word learning. Instead, children get vastly more efficient over
development.

This change in efficiency could come from a variety of sources,
including general maturation of memory and attention, specific changes
in word recognition or processing (\textbf{fernald1998?};
\textbf{frank2026?}), the increasing sophistication of inferential
strategies for word learning
{[}(\textbf{smith2002?});(\textbf{markman1988?});etc{]}, or other
sources.

\begin{figure}[H]

{\centering \pandocbounded{\includegraphics[keepaspectratio]{../figs/longitudinal/exposure_to_learn_EN.png}}

}

\caption{Figure 2.}

\end{figure}%

As an illustration of this change in efficiency, Figure 2 shows a
model-derived estimate of the average number of times a child would need
to hear a word before producing it (See SI: Changes in Word Efficiency),
using frequency estimates from the CHILDES Corpus
(\textbf{mcwhinney2020?}; \textbf{sanchez2019?}). Across different word
classes (e.g., nouns, verbs, closed class words), efficiency increases
exponentially, such that early-learned words such as ``ball'' may be
heard tens of thousands of times before they are produced, while later
learned words such as ``vitamin'' might only be heard tens of times
before being produced.

\subsection{Estimating input}\label{estimating-input}

Under a range of plausible assumptions about variation in input, all
children show supra-linear scaling of growth, even accounting for
accumulation dynamics.

As discussed above, one important question is how much of variance is
due to variance in rate. We explore this by asking about the parameter
\(\pi\),

\[\pi = \frac{\sigma_\alpha^2}{\sigma_\alpha^2 + \sigma_r^2}\]

If \(\pi\) \textgreater{} .5, then the majority of variation between
individuals is not due to pure accumulation rate.

In addition, 80-90\% of between-child variation remains unexplained by
input quantity. These conclusions are supported by analysis of
longitudinal input data from BabyView (N=20) and SEEDLings (N=44)
datasets, which contain child-specific estimates of input, and Peekbank
(N=XYZ), which contains child-specific estimates of processing speed.

Thus, modeling human learning requires positing other sources of
variability and acceleration. Sources of variability could include from
variation in interactional input quality, genetic variation between
children, environmental richness not captured by input, etc. Sources of
acceleration could include increases in processing efficiency,
learning-to-learn including mutual exclusivity, shape bias, syntactic
bootstrapping, as well as domain general changes in perception,
processing speed, memory, attention or other aspects of cognition. The
key to understanding children's data efficiency may thus lie in better
understanding when and how children's learning accelerates with
maturation and experience.

We end by providing a quantitative comparison between within-child human
learning and LLM scaling laws, showing that children's early language
learning is strikingly different from LLM scaling: it shows both steeper
scaling and greater variance in growth rate.

\section{Discussion}\label{discussion}

Coffey 2026 variance decomposition

Blingualism

International adoption

https://bergelsonlab.fas.harvard.edu/sites/g/files/omnuum3941/files/2025-01/meylan\_bergelson\_2021\_ARL.pdf

Claude Code was used for literature review, code generation,
visualization, and simulation infrastructure. All writing is original to
the author.




\end{document}
